\chapterimage{figs/chapterhead.png}
\chapter{The First Model}
\label{chap:first-model}

This chapter is grounded in the current \texttt{WorkInProgress} branch at
commit \texttt{f6dbc298\allowbreak{}66dba3e8\allowbreak{}a7021277\allowbreak{}229b82e9\allowbreak{}6fb54d52}.

The worked example comes from the checked-in model file
\path{models/AirportSecurityExample.gen}.
Its generator is the model-specific source file
\texttt{AirportSecurityExample.cpp}.
Chapter~\ref{chap:first-model} is organized around
Section~\ref{sec:first-model-why-this-example},
Section~\ref{sec:first-model-opening},
Section~\ref{sec:first-model-initial-properties},
Section~\ref{sec:first-model-first-execution},
Section~\ref{sec:first-model-controlled-change}, and
Section~\ref{sec:first-model-what-it-proves}.

\noindent\textbf{Objective.} Walk through a small but complete model from
loading to a first comparison run, using only elements that already exist in
the repository.

\noindent\textbf{Audience.} Readers who have launched the GUI, opened a model
before, and now need one reproducible path that can be inspected, executed, and
slightly modified.

\noindent\textbf{Scope.} Cover the source of the tutorial model, its main
components, the initial simulation settings, the first execution, a single
controlled change, and the limits of what the example proves.

The model choice is intentionally conservative. It uses the familiar flow of
arrival, processing, decision, and disposal. It does not depend on optional
subsystems, and it already exists as a repository artifact. That makes it a
useful bridge between the GUI overview chapter and the later chapters on model
editing and simulation execution.

The chapter follows this sequence:
Figure~\ref{fig:first-model-conceptual-flow},
Figure~\ref{fig:first-model-gui-load},
Figure~\ref{fig:first-model-properties},
Figure~\ref{fig:first-model-simulation-config},
Figure~\ref{fig:first-model-execution},
Figure~\ref{fig:first-model-trace},
Figure~\ref{fig:first-model-report}, and
Figure~\ref{fig:first-model-comparison}.

\section{Why this example}
\label{sec:first-model-why-this-example}

\path{AirportSecurityExample.gen} is a good first model for four reasons.

\begin{itemize}
\item It exists in the repository as a checked-in file.
\item Its structure is small enough to read in one session.
\item It uses the basic modeling primitives that appear throughout the rest of
the user manual.
\item It produces a clear before/after comparison when one parameter is changed.
\end{itemize}

The generated model contains these main elements:
\begin{itemize}
\item \texttt{Create\_1} for passenger arrivals;
\item \texttt{Process\_1} for the security check;
\item \texttt{Decide\_1} for the pass/fail branch;
\item \texttt{Dispose\_1} for cleared passengers;
\item \texttt{Dispose\_2} for denied passengers;
\item \texttt{Transportation\_Security\_Officer} as the seized resource;
\item \texttt{Check\_for\_Proper\_Identification.Queue} as the waiting queue.
\end{itemize}

The source file that generates the model, shown above, confirms that the
example is not invented for the manual. The model is built in code, saved to
\path{models/AirportSecurityExample.gen}, and then executed by the model-
specific application.

\begin{figure}[htbp]
\centering
% FIGURE-SPEC-BEGIN
% ID: fig:first-model-conceptual-flow
% Source of truth:
%   - models/AirportSecurityExample.gen
%   - source/applications/modelSpecific/arenaExamples/AirportSecurityExample.cpp
%   - source/plugins/components/Logic/Create.h
%   - source/plugins/components/DiscreteProcessing/Process.h
%   - source/plugins/components/Decisions/Decide.h
%   - source/plugins/components/Logic/Dispose.h
% Purpose:
%   Show the logical flow of the first tutorial model from arrivals to the two disposal branches.
% Required visual content:
%   - Create -> Process -> Decide chain
%   - resource seize point
%   - queue and branch outcomes
%   - the single simple change point used later in the chapter
% Accessibility description:
%   A compact flow diagram that lets the reader map the textual tutorial to the model structure.
% Validation:
%   The diagram must match the current checked-in example model and its saved component names.
% FIGURE-SPEC-END
\figureplaceholder{figs/generated/first-model-conceptual-flow}
\caption{Conceptual flow of the first tutorial model.}
\label{fig:first-model-conceptual-flow}
\end{figure}

Figure~\ref{fig:first-model-conceptual-flow} is the reference view for the rest
of the chapter. Every later property, run setting, and output should be read
against this flow.

\section{Opening the model in the GUI}
\label{sec:first-model-opening}

Open \path{models/AirportSecurityExample.gen} in the main GUI. After the file
loads, the component tree should expose the model objects listed above, and the
scene should show a short linear process with one branch to each disposal.

The point of this first load is not to edit anything yet. It is to confirm that
the repository model and the GUI agree on the same structure. If the loaded
scene does not resemble the conceptual flow in
Figure~\ref{fig:first-model-conceptual-flow}, stop and recheck the file being
opened.

\begin{figure}[htbp]
\centering
% FIGURE-SPEC-BEGIN
% ID: fig:first-model-gui-load
% Source of truth:
%   - models/AirportSecurityExample.gen
%   - source/applications/gui/genesys/mainwindow.ui
%   - source/applications/gui/genesys/mainwindow.cpp
%   - source/applications/gui/genesys/controllers/ModelInspectorController.cpp
% Purpose:
%   Show the loaded example model inside the current GUI workspace.
% Required visual content:
%   - the model canvas with the short security flow
%   - the component tree
%   - the selected model object
%   - the property editor or nearby inspector
% Accessibility description:
%   A screenshot that proves the tutorial model is actually loaded into the GUI.
% Validation:
%   The scene must show the current repository example, not an unrelated model.
% FIGURE-SPEC-END
\figureplaceholder{figs/generated/first-model-gui-load}
\caption{Loaded tutorial model in the main GUI.}
\label{fig:first-model-gui-load}
\end{figure}

The loaded model should be readable without zoom tricks. The main flow starts at
\texttt{Create\_1}, reaches \texttt{Process\_1}, and then reaches
\texttt{Decide\_1}. The two terminal branches represent the cleared and denied
paths.

\section{Initial properties and simulation settings}
\label{sec:first-model-initial-properties}

The initial model properties are the ones saved in the repository example.
They are the values the reader should see after loading the file for the first
time:

\begin{itemize}
\item arrival distribution: \texttt{expo(2)} minutes;
\item security delay: \texttt{tria(0.75, 1.5, 3)} minutes;
\item pass condition: \texttt{unif(0, 1) < 0.96};
\item replication length: 12 hours;
\item number of replications: 2;
\item warm-up period: 0.5 hour;
\item trace level: event level.
\end{itemize}

These are the properties that make the example reproducible as a tutorial. The
manual does not invent a separate seed field for this model file. Instead, the
reproducible record for a run must include the exact random seed used by the
current execution environment, alongside the file name and the simulation
settings above.

\begin{figure}[htbp]
\centering
% FIGURE-SPEC-BEGIN
% ID: fig:first-model-properties
% Source of truth:
%   - models/AirportSecurityExample.gen
%   - source/applications/gui/genesys/controllers/PropertyEditorController.cpp
%   - source/kernel/simulator/model/ModelSimulation.cpp
% Purpose:
%   Show the key object properties and the simulation values that define the baseline run.
% Required visual content:
%   - property editor with Create_1 selected
%   - the arrival expression and time unit
%   - the process delay and resource request
%   - the simulation configuration summary
% Accessibility description:
%   A close-up of the editable fields that control the tutorial model.
% Validation:
%   The visible properties must match the values stored in AirportSecurityExample.gen.
% FIGURE-SPEC-END
\figureplaceholder{figs/generated/first-model-properties}
\caption{Baseline properties for the tutorial model.}
\label{fig:first-model-properties}
\end{figure}

\begin{figure}[htbp]
\centering
% FIGURE-SPEC-BEGIN
% ID: fig:first-model-simulation-config
% Source of truth:
%   - source/applications/gui/genesys/dialogs/dialogsimulationconfigure.cpp
%   - source/kernel/simulator/model/ModelSimulation.cpp
%   - models/AirportSecurityExample.gen
% Purpose:
%   Show the simulation dialog used to confirm the replication length, warm-up, and replication count.
% Required visual content:
%   - number of replications
%   - replication length
%   - warm-up period
%   - time-unit selectors
% Accessibility description:
%   The simulation settings that make the example run reproducible.
% Validation:
%   The dialog state must reflect the same values stored in the model file.
% FIGURE-SPEC-END
\figureplaceholder{figs/generated/first-model-simulation-config}
\caption{Simulation configuration for the baseline tutorial run.}
\label{fig:first-model-simulation-config}
\end{figure}

Before the first execution, confirm two things. First, the model file name is
the one above, not a different saved copy. Second, the simulation settings in
the dialog match the saved model values. If either one is different, the run is
no longer the baseline version described in this chapter.

\section{First execution}
\label{sec:first-model-first-execution}

Run the model once without changing anything. The expected behavior is simple:
passengers arrive, enter the process, seize the officer resource, wait in the
queue if needed, and leave through one of the two disposal paths.

Because the tutorial uses stochastic distributions, the exact event order and
the exact counts depend on the seed used in the run. That does not weaken the
tutorial. It only means that the reproducible record must keep the seed with
the model file and the simulation settings.

\begin{figure}[htbp]
\centering
% FIGURE-SPEC-BEGIN
% ID: fig:first-model-execution
% Source of truth:
%   - source/applications/gui/genesys/mainwindow.cpp
%   - source/kernel/simulator/model/ModelSimulation.cpp
%   - models/AirportSecurityExample.gen
% Purpose:
%   Show the GUI while the baseline simulation is actively running.
% Required visual content:
%   - active run indicator
%   - current simulation time or progress
%   - visible model or trace updates
%   - controls in a state consistent with execution
% Accessibility description:
%   A screenshot that proves the simulation is running and not just configured.
% Validation:
%   The running state must correspond to the baseline example described in the text.
% FIGURE-SPEC-END
\figureplaceholder{figs/generated/first-model-execution}
\caption{Baseline simulation running in the GUI.}
\label{fig:first-model-execution}
\end{figure}

At this stage, the main thing to verify is not a specific number but a correct
qualitative path: arrivals happen, the security step consumes time, and the
branching decision sends some entities to the cleared path and some to the
denied path.

\begin{figure}[htbp]
\centering
% FIGURE-SPEC-BEGIN
% ID: fig:first-model-trace
% Source of truth:
%   - source/applications/gui/genesys/mainwindow.ui
%   - source/kernel/simulator/model/ModelSimulation.cpp
%   - source/kernel/trace/TraceManager.h
% Purpose:
%   Show the trace or event log produced by the first run.
% Required visual content:
%   - ordered event messages
%   - at least one arrival, one processing step, and one branch result
%   - the event-level trace setting used for the example
% Accessibility description:
%   A readable trace view that connects the flow diagram to runtime behavior.
% Validation:
%   The trace must match the baseline example and show real runtime messages.
% FIGURE-SPEC-END
\figureplaceholder{figs/generated/first-model-trace}
\caption{Trace view for the first baseline execution.}
\label{fig:first-model-trace}
\end{figure}

\begin{figure}[htbp]
\centering
% FIGURE-SPEC-BEGIN
% ID: fig:first-model-report
% Source of truth:
%   - source/kernel/simulator/model/ModelSimulation.cpp
%   - source/applications/gui/genesys/mainwindow.ui
%   - models/AirportSecurityExample.gen
% Purpose:
%   Show the output report or summary statistics generated by the baseline run.
% Required visual content:
%   - reported replication summary
%   - any observed statistics for the branch outcomes
%   - time-based summary values if visible
% Accessibility description:
%   A report snapshot that lets the reader relate the run to the observed statistics.
% Validation:
%   The report must reflect the same example and the same simulation settings as the trace.
% FIGURE-SPEC-END
\figureplaceholder{figs/generated/first-model-report}
\caption{Report summary from the baseline execution.}
\label{fig:first-model-report}
\end{figure}

The trace and the report together are enough to establish that the example is
working. A successful first run proves that the GUI can load the model, the
kernel can execute it, and the reporting path can summarize it.

\section{One controlled change}
\label{sec:first-model-controlled-change}

The simplest useful change is to make the security branch stricter. Change the
decision condition from \texttt{unif(0, 1) < 0.96} to
\texttt{unif(0, 1) < 0.90}. That keeps the structure intact but increases the
number of entities that fall into the denied path.

The change is deliberately small:
\begin{itemize}
\item no new components are added;
\item no connections are rewired;
\item only one numeric threshold changes;
\item the effect should be visible in the trace and report.
\end{itemize}

Save the modified model under a new name before rerunning it. That preserves
the baseline example and makes the comparison explicit.

\begin{figure}[htbp]
\centering
% FIGURE-SPEC-BEGIN
% ID: fig:first-model-comparison
% Source of truth:
%   - models/AirportSecurityExample.gen
%   - the saved modified copy created during the tutorial
%   - source/kernel/simulator/model/ModelSimulation.cpp
% Purpose:
%   Show the before/after comparison between the baseline run and the stricter decision threshold.
% Required visual content:
%   - baseline and modified report summary side by side, or a clearly comparable pair
%   - the changed decision threshold in the edited model
%   - a visible difference in branch counts or branch proportions
% Accessibility description:
%   A comparison figure that makes the single parameter change and its effect obvious.
% Validation:
%   The comparison must show a real difference caused by the edited threshold, not by a different model.
% FIGURE-SPEC-END
\figureplaceholder{figs/generated/first-model-comparison}
\caption{Comparison between the baseline and the stricter branch threshold.}
\label{fig:first-model-comparison}
\end{figure}

The comparison should answer one question clearly: did the changed threshold
affect the outcome in the expected direction? If the answer is yes, the user
has completed a full first-model loop:
load, read, run, change, rerun, compare.

\section{What this example proves}
\label{sec:first-model-what-it-proves}

This tutorial proves that the checked-in example model can be loaded and
understood as a small simulation system. It also proves that a single
parameter change can produce an observable difference in the result path.

What it does not prove is broader than that:
\begin{itemize}
\item it does not certify the statistical quality of the chosen distributions;
\item it does not establish a universal seed policy for all models;
\item it does not replace the later chapters on editing, execution, and results;
\item it does not claim the example is a final product model.
\end{itemize}

That limitation is intentional. A first model should be concrete, readable, and
repeatable, not exhaustive. The later chapters build on the same logic with
more detailed workflows and richer example structures.
